\documentclass[UTF8,12pt]{ctexart}
\usepackage[a4paper,margin=24mm]{geometry}
\usepackage{amsmath,amssymb,bm,graphicx,adjustbox}
\newcommand{\fitmath}[1]{\begin{adjustbox}{max width=\linewidth}$\displaystyle #1$\end{adjustbox}}
\setlength{\parindent}{0pt}
\setlength{\parskip}{.7em}
\sloppy
\allowdisplaybreaks
\begin{document}
\section*{1992 年考研数学三 · 真题}
\subsection*{1992年全国硕士研究生招生考试试题}

\subsection*{（试卷IV）}

\subsection*{一、填空题（本题共5小题，每小题3分，满分15分）}

（1）设商品的需求函数为 \(\displaystyle Q=100-5p\)，其中 \(\displaystyle Q,\allowbreak p\) 分别表示需求量和价格，如果商品需求弹性的绝对值大于1，则商品价格的取值范围是\_\_\_\_\_\_。


（2）级数 \(\displaystyle \sum_{n=1}^{\infty}\dfrac{\displaystyle (x-2)^{2n}}{\displaystyle n4^n}\) 的收敛域为\_\_\_\_\_\_。


（3）交换积分次序 \(\displaystyle \int_0^1dy\displaystyle \int_{\sqrt y}^{\sqrt{2-y^2}}f(x,\allowbreak y)\,dx=\)\_\_\_\_\_\_。


（4）设 \(\displaystyle A\) 为 \(\displaystyle m\) 阶方阵，\(\displaystyle B\) 为 \(\displaystyle n\) 阶方阵，且 \(\displaystyle |A|=a,\allowbreak |B|=b\)，\(\displaystyle C=\begin{pmatrix}O&A\\B&O\end{pmatrix}\)，则 \(\displaystyle |C|=\)\_\_\_\_\_\_。


（5）将 \(\displaystyle C,\allowbreak C,\allowbreak E,\allowbreak E,\allowbreak I,\allowbreak N,\allowbreak S\) 等七个字母随机地排成一行，那么，恰好排成 SCIENCE 的概率为\_\_\_\_\_\_。


\subsection*{二、选择题（本题共5小题，每小题3分，满分15分）}

（1）设 \(\displaystyle F(x)=\dfrac{\displaystyle x^2}{\displaystyle x-a}\displaystyle \int_a^x f(t)\,dt\)，其中 \(\displaystyle f(x)\) 为连续函数，则 \(\displaystyle \lim_{x\to a}F(x)\) 等于（　　） （A）\(\displaystyle a^2\)。 （B）\(\displaystyle a^2f(a)\)。 （C）\(\displaystyle 0\)。 （D）不存在。


（2）当 \(\displaystyle x\to0\) 时，下列四个无穷小量中，哪一个是比其他三个更高阶的无穷小量？（　　） （A）\(\displaystyle x^2\)。 （B）\(\displaystyle 1-\cos x\)。 （C）\(\displaystyle \sqrt{1-x^2}-1\)。 （D）\(\displaystyle x-\tan x\)。


（3）设 \(\displaystyle A\) 为 \(\displaystyle m\times n\) 矩阵，则齐次线性方程组 \(\displaystyle A\boldsymbol x=0\) 仅有零解的充分条件是（　　） （A）\(\displaystyle A\) 的列向量线性无关。 （B）\(\displaystyle A\) 的列向量线性相关。 （C）\(\displaystyle A\) 的行向量线性无关。 （D）\(\displaystyle A\) 的行向量线性相关。


（4）设当事件 \(\displaystyle A\) 与 \(\displaystyle B\) 同时发生时，事件 \(\displaystyle C\) 必发生，则（　　） （A）\(\displaystyle P(C)\le P(A)+P(B)-1\)。 （B）\(\displaystyle P(C)\ge P(A)+P(B)-1\)。 （C）\(\displaystyle P(C)=P(AB)\)。 （D）\(\displaystyle P(C)=P(A\cup B)\)。


（5）（超纲题）设 \(\displaystyle n\) 个随机变量 \(\displaystyle X_1,\allowbreak X_2,\allowbreak \ldots,\allowbreak X_n\) 独立同分布，\(\displaystyle D(X_1)=\sigma^2\)，\(\displaystyle \bar X=\dfrac{\displaystyle 1}{\displaystyle n}\displaystyle \sum_{i=1}^nX_i\)，\(\displaystyle S^2=\dfrac{\displaystyle 1}{\displaystyle n-1}\displaystyle \sum_{i=1}^n(X_i-\bar X)^2\)，则（　　）


（A）\(\displaystyle S\) 是 \(\displaystyle \sigma\) 的无偏估计量（可改为“\(\displaystyle E(S)=\sigma\)”）。 （B）\(\displaystyle S\) 是 \(\displaystyle \sigma\) 的最大似然估计量。 （C）\(\displaystyle S\) 是 \(\displaystyle \sigma\) 的相合估计量（即一致估计量）。 （D）\(\displaystyle S\) 与 \(\displaystyle \bar X\) 相互独立。


\subsection*{三、（本题满分5分）}

设函数


\[\fitmath{\displaystyle f(x)=\begin{cases}\dfrac{\displaystyle \ln[\cos(x-1)]}{\displaystyle 1-\sin(\dfrac{\displaystyle \pi}{\displaystyle 2}x)},&x\ne1,\\1,&x=1,\end{cases}}\]


问函数 \(\displaystyle f(x)\) 在 \(\displaystyle x=1\) 处是否连续？若不连续，修改函数在 \(\displaystyle x=1\) 处的定义使之连续。


\subsection*{四、（本题满分5分）}

计算 \(\displaystyle I=\displaystyle \int\dfrac{\displaystyle \operatorname{arccot}e^x}{\displaystyle e^x}\,dx\)。


\subsection*{五、（本题满分5分）}

设 \(\displaystyle z=\sin(xy)+\varphi(x,\allowbreak \dfrac{\displaystyle x}{\displaystyle y})\)，求 \(\displaystyle \dfrac{\displaystyle \partial^2z}{\displaystyle \partial x\partial y}\)，其中 \(\displaystyle \varphi(u,\allowbreak v)\) 有二阶偏导数。


\subsection*{六、（本题满分5分）}

求连续函数 \(\displaystyle f(x)\)，使它满足 \(\displaystyle f(x)+2\displaystyle \int_0^x f(t)\,dt=x^2\)。


\subsection*{七、（本题满分6分）}

求证：当 \(\displaystyle x\ge1\) 时，\(\displaystyle \arctan x-\dfrac{\displaystyle 1}{\displaystyle 2}\arccos\dfrac{\displaystyle 2x}{\displaystyle 1+x^2}=\dfrac{\displaystyle \pi}{\displaystyle 4}\)。


\subsection*{八、（本题满分9分）}

设曲线方程为 \(\displaystyle y=e^{-x}\;(x\ge0)\)。（1）把曲线 \(\displaystyle y=e^{-x}\)、\(\displaystyle x\) 轴、\(\displaystyle y\) 轴和直线 \(\displaystyle x=\xi\;(\xi>0)\) 所围平面图形绕 \(\displaystyle x\) 轴旋转一周，得一旋转体，求此旋转体体积 \(\displaystyle V(\xi)\)；求满足 \(\displaystyle V(a)=\dfrac{\displaystyle 1}{\displaystyle 2}\displaystyle \lim_{\xi\to+\infty}V(\xi)\) 的 \(\displaystyle a\)。（2）在此曲线上找一点，使过该点的切线与两个坐标轴所夹平面图形的面积最大，并求出该面积。


\subsection*{九、（本题满分7分）}

设矩阵 \(\displaystyle A\) 与 \(\displaystyle B\) 相似，其中


\[\fitmath{\displaystyle A=\begin{pmatrix}-2&0&0\\2&x&2\\3&1&1\end{pmatrix},\quad B=\begin{pmatrix}-1&0&0\\0&2&0\\0&0&y\end{pmatrix}.}\]


（1）求 \(\displaystyle x\) 和 \(\displaystyle y\) 的值；（2）求可逆矩阵 \(\displaystyle P\)，使 \(\displaystyle P^{-1}AP=B\)。


\subsection*{十、（本题满分6分）}

已知3阶矩阵 \(\displaystyle B\ne O\)，且 \(\displaystyle B\) 的每一个列向量都是以下方程组的解：


\[\fitmath{\displaystyle \begin{cases}x_1+2x_2-2x_3=0,\\2x_1-x_2+\lambda x_3=0,\\3x_1+x_2-x_3=0.\end{cases}}\]


（1）求 \(\displaystyle \lambda\) 的值；（2）证明 \(\displaystyle |B|=0\)。


\subsection*{十一、（本题满分6分）}

设 \(\displaystyle A,\allowbreak B\) 分别为 \(\displaystyle m\) 阶、\(\displaystyle n\) 阶正定矩阵，试判定分块矩阵 \(\displaystyle C=\begin{pmatrix}A&O\\O&B\end{pmatrix}\) 是否是正定矩阵。


\subsection*{十二、（本题满分7分）}

假设测量的随机误差 \(\displaystyle X\sim N(0,\allowbreak 10^2)\)，试求在100次独立重复测量中，至少有三次测量误差的绝对值大于19.6的概率 \(\displaystyle \alpha\)，并利用泊松分布求出 \(\displaystyle \alpha\) 的近似值（要求小数点后取两位有效数字）。附表：


\[\fitmath{\displaystyle \begin{array}{c|cccccccc}\lambda&1&2&3&4&5&6&7&\cdots\\\hline e^{-\lambda}&0.368&0.135&0.050&0.018&0.007&0.002&0.001&\cdots\end{array}}\]


\subsection*{十三、（本题满分5分）}

一台设备由三大部件构成，在设备运转中各部件需要调整的概率相应为0.10、0.20和0.30。假设各部件的状态相互独立，以 \(\displaystyle X\) 表示同时需要调整的部件数，试求 \(\displaystyle X\) 的概率分布、数学期望 \(\displaystyle E(X)\) 和方差 \(\displaystyle D(X)\)。


\subsection*{十四、（本题满分4分）}

设二维随机变量 \(\displaystyle (X,\allowbreak Y)\) 的概率密度为


\[\fitmath{\displaystyle f(x,y)=\begin{cases}e^{-y},&0<x<y,\\0,&\text{其他},\end{cases}}\]


（1）求 \(\displaystyle X\) 的概率密度 \(\displaystyle f_X(x)\)；（2）求概率 \(\displaystyle P\{X+Y\le1\}\)。


\subsection*{（试卷V）}

\subsection*{一、填空题（本题共5小题，每小题3分，满分15分）}

（1）设 \(\displaystyle f(t)=\displaystyle \lim_{x\to\infty}t\left(\dfrac{\displaystyle x+t}{\displaystyle x-t}\right)^x\)，则 \(\displaystyle f\prime(t)=\)\_\_\_\_\_\_。


（2）【同试卷IV 第一、（1）题】


（3）设 \(\displaystyle f(x)=\sin x\)，\(\displaystyle f[\varphi(x)]=1-x^2\)，则 \(\displaystyle \varphi(x)=\)\_\_\_\_\_\_；其定义域为\_\_\_\_\_\_。


（4）矩阵


\[\fitmath{\displaystyle A=\begin{pmatrix}1&1&1&1\\1&1&1&1\\1&1&1&1\\1&1&1&1\end{pmatrix}}\]


的非零特征值是\_\_\_\_\_\_。


（5）设对于事件 \(\displaystyle A,\allowbreak B,\allowbreak C\)，有 \(\displaystyle P(A)=P(B)=P(C)=\dfrac{\displaystyle 1}{\displaystyle 4},\allowbreak \ P(AB)=P(BC)=0,\allowbreak \ P(AC)=\dfrac{\displaystyle 1}{\displaystyle 8}\)，则 \(\displaystyle A,\allowbreak B,\allowbreak C\) 三个事件中至少出现一个的概率为\_\_\_\_\_\_。


\subsection*{二、选择题（本题共5小题，每小题3分，满分15分）}

（1）【同试卷IV 第二、（1）题】


（2）当 \(\displaystyle x\to0\) 时，下列四个无穷小量中，比其它三个更高阶的无穷小量是（　　） （A）\(\displaystyle x^2\)。 （B）\(\displaystyle 1-\cos x\)。 （C）\(\displaystyle \sqrt{1-x^2}-1\)。 （D）\(\displaystyle x-\sin x\)。


（3）设 \(\displaystyle A,\allowbreak B,\allowbreak A+B,\allowbreak A^{-1}+B^{-1}\) 均为 \(\displaystyle n\) 阶可逆矩阵，则 \(\displaystyle (A^{-1}+B^{-1})^{-1}\) 等于（　　） （A）\(\displaystyle A^{-1}+B^{-1}\)。 （B）\(\displaystyle A+B\)。 （C）\(\displaystyle A(A+B)^{-1}B\)。 （D）\(\displaystyle (A+B)^{-1}\)。


（4）设 \(\displaystyle \boldsymbol\alpha_1,\allowbreak \boldsymbol\alpha_2,\allowbreak \ldots,\allowbreak \boldsymbol\alpha_m\) 均为 \(\displaystyle n\) 维向量，那么下列结论正确的是（　　）


（A）若 \(\displaystyle k_1\boldsymbol\alpha_1+\cdots+k_m\boldsymbol\alpha_m=0\)，则 \(\displaystyle \boldsymbol\alpha_1,\allowbreak \ldots,\allowbreak \boldsymbol\alpha_m\) 线性相关。 （B）若对任意一组不全为零的数 \(\displaystyle k_1,\allowbreak \ldots,\allowbreak k_m\)，都有 \(\displaystyle k_1\boldsymbol\alpha_1+\cdots+k_m\boldsymbol\alpha_m\ne0\)，则 \(\displaystyle \boldsymbol\alpha_1,\allowbreak \ldots,\allowbreak \boldsymbol\alpha_m\) 线性无关。


（C）若 \(\displaystyle \boldsymbol\alpha_1,\allowbreak \ldots,\allowbreak \boldsymbol\alpha_m\) 线性相关，则对任意一组不全为零的数 \(\displaystyle k_1,\allowbreak \ldots,\allowbreak k_m\)，都有 \(\displaystyle k_1\boldsymbol\alpha_1+\cdots+k_m\boldsymbol\alpha_m=0\)。 （D）若 \(\displaystyle 0\boldsymbol\alpha_1+\cdots+0\boldsymbol\alpha_m=0\)，则 \(\displaystyle \boldsymbol\alpha_1,\allowbreak \ldots,\allowbreak \boldsymbol\alpha_m\) 线性无关。


（5）【同试卷IV 第二、（4）题】


\subsection*{三、（本题满分5分）}

求极限 \(\displaystyle \lim_{x\to1}\dfrac{\displaystyle \ln[\cos(x-1)]}{\displaystyle 1-\sin(\dfrac{\displaystyle \pi}{\displaystyle 2}x)}\)。


\subsection*{四、（本题满分5分）【同试卷IV 第四题】}

\subsection*{五、（本题满分6分）}

求连续函数 \(\displaystyle f(x)\)，使它满足 \(\displaystyle \int_0^1f(tx)\,dt=f(x)+x\sin x\)。


\subsection*{六、（本题满分5分）【同试卷IV 第五题】}

\subsection*{七、（本题满分6分）}

设生产某产品的固定成本为10，而当产量为 \(\displaystyle x\) 时的边际成本函数为 \(\displaystyle MC=-40-20x+3x^2\)，边际收入函数为 \(\displaystyle MR=32+10x\)。试求：（1）总利润函数；（2）使总利润最大的产量。


\subsection*{八、（本题满分6分）}

求证：方程 \(\displaystyle x+p+q\cos x=0\) 恰有一个实根，其中 \(\displaystyle p,\allowbreak q\) 为常数，且 \(\displaystyle 0<q<1\)。


\subsection*{九、（本题满分7分）}

给定曲线 \(\displaystyle y=\dfrac{\displaystyle 1}{\displaystyle x^2}\)。（1）求曲线在横坐标为 \(\displaystyle x_0\) 的点处的切线方程；


（2）求曲线的切线被两坐标轴所截线段的最短长度。


\subsection*{十、（本题满分5分）}

设矩阵 \(\displaystyle X\) 满足 \(\displaystyle AX+E=A^2+X\)，其中 \(\displaystyle E\) 为3阶单位阵，又已知 \(\displaystyle A=\begin{pmatrix}1&0&1\\0&2&0\\1&0&1\end{pmatrix}\)，试求出矩阵 \(\displaystyle X\)。


\subsection*{十一、（本题满分5分）}

设线性方程组


\[\fitmath{\displaystyle \begin{cases}x_1+2x_2-2x_3=0,\\2x_1-x_2+\lambda x_3=0,\\3x_1+x_2-x_3=0\end{cases}}\]


的系数矩阵为 \(\displaystyle A\)，3阶矩阵 \(\displaystyle B\ne O\)，且 \(\displaystyle AB=O\)。试求 \(\displaystyle \lambda\) 的值。


\subsection*{十二、（本题满分6分）}

已知实矩阵 \(\displaystyle A=(a_{ij})_{3\times3}\) 满足条件：\textcircled{1} \(\displaystyle a_{ij}=A_{ij}\;(i,\allowbreak j=1,\allowbreak 2,\allowbreak 3)\)，其中 \(\displaystyle A_{ij}\) 是 \(\displaystyle a_{ij}\) 的代数余子式；\textcircled{2} \(\displaystyle a_{11}\ne0\)。计算行列式 \(\displaystyle |A|\)。


\subsection*{十三、（本题满分7分）【同试卷IV 第十二题】}

\subsection*{十四、（本题满分7分）【同试卷IV 第三题】}
\end{document}
